\paragraph{Tree realization of a normally generated target.}
Let $X$ be an actual two-dimensional $K(H,1)$, let $q:H\to E$ be a
homomorphism, and suppose that $q(H)$ normally generates $E$. There is an
actual two-dimensional aspherical complex $Y\supseteq X$ such that
$H\hookrightarrow\pi_1Y$ is a homology isomorphism and normally generates
$\pi_1Y$, and $q$ extends to an epimorphism $\pi_1Y\twoheadrightarrow E$.
The number of new cells is at most $\max(\aleph_0,|E|)$. This statement
does not say that a prescribed homomorphism to a subgroup of $E$ must
extend surjectively onto $E$.

\emph{Proof.}
For each $e\in E$ choose a root labelled $e$. For any node $s$ labelled
$e_s$, choose a finite expression
\[
 e_s=\prod_{j=1}^{r_s}u_{sj}\,q(h_{sj})\,u_{sj}^{-1},
 \qquad h_{sj}\in H,
\]
and give it $r_s$ new, distinct children labelled $u_{sj}$.
Repeat at all nodes, using distinct variables even when labels repeat.
Attach a loop $x_s$ and a two-cell with boundary
\[
             x_s^{-1}\prod_{j=1}^{r_s}
                    x_{sj}h_{sj}x_{sj}^{-1}
\]
for every node $s$. Empty products are permitted.
The assignment $x_s\mapsto e_s$ extends $q$, and the root labels make
the resulting homomorphism onto $E$.

Every finite set of cells is contained in the complex consisting of
$X$, finitely many finite tree prefixes, and the loops at their
frontiers. Starting at the roots, each internal loop occurs once in
its own two-cell and in no other two-cell remaining after its ancestors
have been eliminated. Cancelling these generator--relator pairs
gives a homotopy equivalence to $X$ wedged with the frontier circles.
Such a finite prefix is aspherical and contains $H$ injectively.
Compactness of spheres and disk homotopies now gives asphericity of
the actual union $Y$ and injectivity of $H\to\pi_1Y$; no assumption
that successive frontier groups inject is needed.

The relative cellular differential $C_2(Y,X)\to C_1(Y,X)$ sends the
cell at $s$ to $-x_s$: each child has exponent $+1-1=0$, while every
coefficient lies in $X$. It is therefore the negative identity on
the free abelian groups indexed by all nodes. Hence
$H_*(X;\mathbb Z)\to H_*(Y;\mathbb Z)$ is an isomorphism. In the
quotient of $\pi_1Y$ by the normal closure of $H$, each defining
relation directly gives $x_s=1$, since its right side is a product
of conjugates of trivial coefficients. Thus $H$ normally generates.
Finite branching at each of countably many levels gives the stated
cell bound.\hfill$\square$

\emph{Why this does not transfer a killer from $E$ to an image subgroup.}
The map to $E$ constructed above is only one extension of $q$.
Another homomorphism out of $\pi_1Y$ can assign different values to the
tree variables. In particular, if all the chosen coefficient products
$\prod_jh_{sj}$ equal $1$ in $H$, sending every $x_s$ to $1$ gives a
retraction $\pi_1Y\to H$. It is therefore invalid to apply a later
obstruction to this one epimorphism onto $E$ and conclude that all
extensions into $q(H)$ have been obstructed. The construction supplies
coefficient lifts for its selected map, not uniqueness or forced
surjectivity for arbitrary maps.
